(1) 记 $x=\Delta/(k_BT)$，配分函数 $Z=1+e^{-x}$，高能级概率 $p_1=e^{-x}/Z=1/(1+e^x)$。平均能量 $U=\Delta p_1=\Delta/(1+e^x)$，随温度上升从零逐渐趋近 $\Delta/2$。
(2) 因 $dx/dT=-x/T$，定容热容 $C=dU/dT=k_Bx^2e^x/(1+e^x)^2=k_Bx^2/[4\cosh^2(x/2)]$。低温 $x\gg1$ 时 $C\sim k_Bx^2e^{-x}$ 指数小；高温 $x\ll1$ 时 $C\sim k_Bx^2/4=k_B\Delta^2/(4k_B^2T^2)$，均趋零，因此中间必有峰。
(3) 对 $\ln(C/k_B)=2\ln x-2\ln[2\cosh(x/2)]$ 求导，极值方程 $2/x=\tanh(x/2)$。对 $x>0$，左侧严格下降，右侧严格上升，从而唯一交点 $x_*\simeq2.39936$，峰温 $T_*=\Delta/(k_Bx_*)$；峰值有限，是离散能隙被热激发的 Schottky 异常。
(4) 由正则熵公式 $S=k_B(\ln Z+\beta U)$，有 $S=k_B[\ln(1+e^{-x})+x/(1+e^x)]$。$T\to0$ 时 $x\to\infty$ 且 $S\to0$，因为基态唯一；$T\to\infty$ 时 $x\to0$，$S\to k_B\ln2$。有限两态配分函数在所有 $T>0$ 解析，故不能把有限热容峰误认为热力学相变。
